\chapter{Derivative Estimates}

\section*{1. Global estimates and their consequences}

\begin{theorem}[Global derivative estimates]
\label{ch07-thm-global-derivative-estimates}
\source{chow-knopf-ricci-flow:page-0236, chow-knopf-ricci-flow:page-0237, chow-knopf-ricci-flow:page-0240, chow-knopf-ricci-flow:page-0241, chow-knopf-ricci-flow:page-0242, chow-knopf-ricci-flow:page-0243, chow-knopf-ricci-flow:page-0244}
\dcref{Theorem 7.1}
\group{global-estimates-and-their-consequences}

Let \((\mathcal M^n,g(t))\) be a solution of the Ricci flow for which the maximum principle holds. Then for each \(\alpha>0\) and every \(m\in\mathbb N\), there exists a constant \(C_m\) depending only on \(m\), \(n\), and \(\max\{\alpha,1\}\) such that if
\[
|\Rm(x,t)|_{g(x,t)}\le K
\qquad\text{for all }x\in \mathcal M^n\text{ and }t\in\left[0,\frac{\alpha}{K}\right],
\]
then
\[
|\nabla^m \Rm(x,t)|_{g(x,t)}\le \frac{C_m K}{t^{m/2}}
\qquad\text{for all }x\in \mathcal M^n\text{ and }t\in\left(0,\frac{\alpha}{K}\right].
\]
\end{theorem}

\begin{proof}
\source{chow-knopf-ricci-flow:page-0240, chow-knopf-ricci-flow:page-0241, chow-knopf-ricci-flow:page-0242, chow-knopf-ricci-flow:page-0243, chow-knopf-ricci-flow:page-0244}
\sketch The proof is by induction on \(m\). For \(m=1\), the evolution equation for \(|\nabla \Rm|^2\) has the form
\[
\partial_t |\nabla \Rm|^2=\Delta |\nabla \Rm|^2-2|\nabla^2\Rm|^2+\Rm*(\nabla\Rm)^{*2}.
\]
Introducing
\[
F=t|\nabla \Rm|^2+\beta |\Rm|^2
\]
and choosing \(\beta\) large enough in terms of \(n\) and \(\max\{\alpha,1\}\), the maximum principle gives a uniform bound on \(F\), hence
\[
|\nabla \Rm|\le C_1K\,t^{-1/2}.
\]
For the inductive step, one computes the evolution of \(|\nabla^m\Rm|^2\) and uses the commutator identity
\[
[\nabla^m,\Delta]A=\sum_{j=0}^m \nabla^j\Rm*\nabla^{m-j}A
\]
to obtain a differential inequality of the form
\[
\partial_t |\nabla^m\Rm|^2
\le \Delta |\nabla^m\Rm|^2+\bar C_m K\Bigl(|\nabla^m\Rm|^2+\frac{K^2}{t^m}\Bigr).
\]
Defining
\[
G=t^m|\nabla^m\Rm|^2+\beta_m \sum_{k=1}^m \frac{(m-1)!}{(m-k)!}\,t^{m-k}|\nabla^{m-k}\Rm|^2,
\]
one arranges cancellation between the unfavorable terms coming from differentiating the time weights and the favorable negative terms in the lower-order estimates. Choosing \(\beta_m\) sufficiently large, the maximum principle gives a bound \(\sup G\le C_m^2K^2\), and therefore
\[
|\nabla^m\Rm|\le C_mK\,t^{-m/2}.
\]
\end{proof}

\begin{corollary}[Long-time existence criterion]
\label{ch07-cor-long-time-existence}
\source{chow-knopf-ricci-flow:page-0237}
\dcref{Corollary 7.2}
\group{global-estimates-and-their-consequences}

If \(g_0\) is a smooth metric on a compact manifold \(\mathcal M^n\), the unique Ricci flow solution \(g(t)\) with \(g(0)=g_0\) exists on a maximal time interval \(0\le t<T\le \infty\). Moreover, \(T<\infty\) only if
\[
\lim_{t\nearrow T}\left(\sup_{x\in \mathcal M^n}|\Rm(x,t)|\right)=\infty.
\]
\end{corollary}

\begin{proof}
\source{chow-knopf-ricci-flow:page-0237}
\sketch If the curvature remained bounded on a finite maximal interval, then Theorem~\dcref{Theorem 7.1} would control all curvature derivatives on a short time interval near the terminal time. Standard short-time existence then extends the flow past \(T\), contradicting maximality. Hence finite-time breakdown forces curvature blow-up.
\end{proof}

\begin{corollary}[Uniform derivative bounds for bounded-curvature sequences]
\label{ch07-cor-uniform-bounds-sequences}
\source{chow-knopf-ricci-flow:page-0237}
\dcref{Corollary 7.3}
\group{global-estimates-and-their-consequences}

Let \(\{(\mathcal M_i^n,g_i^0):i\in\mathbb N\}\) be a sequence of compact Riemannian manifolds. If their curvatures satisfy
\[
|\Rm[g_i^0]|_{g_i^0}\le K
\]
for some \(K<\infty\) independent of \(i\), then there exists \(c>0\) depending only on \(n\) such that each Ricci flow solution \(g_i(t)\) with \(g_i(0)=g_i^0\) exists on \(t\in[0,c/K]\). Moreover, for each \(m\in\mathbb N\) there exists \(C_m<\infty\) depending only on \(m\) and \(n\) such that
\[
\sup_{x\in \mathcal M_i^n}|\nabla^m\Rm[g_i(t)]|_{g_i(t)}\le \frac{C_mK}{(\sqrt t)^m}.
\]
In particular, for any \(t_0\in(0,c/K)\), there is a constant \(C_m'=C_m'(m,n,K,t_0)\) such that
\[
\sup_{\mathcal M_i^n\times[t_0,c/K]}|\nabla^m\Rm[g_i]|_{g_i}\le C_m'.
\]
\end{corollary}

\begin{proof}
\source{chow-knopf-ricci-flow:page-0237}
\sketch Apply the long-time existence criterion after first using the minimal existence time estimate of Corollary~\dcref{Corollary 7.7}, which gives a uniform time interval \([0,c/K]\) depending only on \(n\). Then Theorem~\dcref{Theorem 7.1} gives the derivative estimates on that interval. The final bound on \([t_0,c/K]\) is immediate because \(t^{-m/2}\) is uniformly bounded there.
\end{proof}

\section*{2. Proving the global estimates}

\begin{lemma}[Evolution of the curvature norm]
\label{ch07-lem-curvature-norm-evolution}
\source{chow-knopf-ricci-flow:page-0238}
\dcref{Lemma 7.4}
\group{proving-the-global-estimates}

If \((\mathcal M^n,g(t))\) is a solution of the Ricci flow, then the square of the norm of its curvature tensor evolves by
\[
\frac{\partial}{\partial t}|\Rm|^2
=\Delta |\Rm|^2-2|\nabla \Rm|^2
+4g^{ri}g^{sj}g^{pk}g^{q\ell}R_{rspq}\bigl(B_{ijk\ell}-B_{ij\ell k}+B_{ikj\ell}-B_{i\ell jk}\bigr),
\]
where
\[
B_{ijk\ell}=-R^q{}_{pij}R^p{}_{qk\ell}.
\]
In particular,
\[
\frac{\partial}{\partial t}|\Rm|^2\le \Delta |\Rm|^2-2|\nabla \Rm|^2+C|\Rm|^3
\]
for a constant \(C\) depending only on \(n\).
\end{lemma}

\begin{proof}
\source{chow-knopf-ricci-flow:page-0238}
\sketch Differentiate the \((4,0)\)-curvature evolution equation from Chapter~6. The terms from differentiating the inverse metric cancel exactly with the Ricci terms in the curvature evolution, leaving the displayed formula. The inequality follows because the full contraction of two curvature tensors is bounded by a dimensional constant times \(|\Rm|^3\).
\end{proof}

\begin{corollary}[Doubling-time estimate]
\label{ch07-cor-doubling-time-estimate}
\source{chow-knopf-ricci-flow:page-0238, chow-knopf-ricci-flow:page-0239}
\dcref{Corollary 7.5}
\group{proving-the-global-estimates}

There exists \(c>0\) depending only on the dimension \(n\) such that if \((\mathcal M^n,g(t):0\le t\le \tau)\) is a solution of the Ricci flow on a compact manifold and
\[
M(t)\coloneqq \sup_{x\in\mathcal M^n}|\Rm(x,t)|_{g(x,t)},
\]
then
\[
M(t)\le 2M(0)
\qquad\text{for all }0\le t\le \min\left\{\tau,\frac{c}{M(0)}\right\}.
\]
\end{corollary}

\begin{proof}
\source{chow-knopf-ricci-flow:page-0239}
\sketch By Lemma~\dcref{Lemma 7.4}, \(M(t)\) is locally Lipschitz and satisfies an upper differential inequality of the form
\[
\frac{dM}{dt}\le \frac{C}{2}M^2
\]
in the forward-difference sense. Solving the corresponding ODE yields
\[
M(t)\le \frac{1}{M(0)^{-1}-\frac C2 t}
\]
while this denominator stays positive. Taking \(c=1/C\) gives \(M(t)\le 2M(0)\) for \(t\le \min\{\tau,c/M(0)\}\).
\end{proof}

\begin{remark}[A special case]
\label{ch07-rem-special-case-doubling-time}
\source{chow-knopf-ricci-flow:page-0239}
\dcref{Remark 7.6}
\group{proving-the-global-estimates}

The preceding doubling-time estimate generalizes the special case noted earlier in the book.
\end{remark}

\begin{corollary}[Minimal existence time]
\label{ch07-cor-minimal-existence-time}
\source{chow-knopf-ricci-flow:page-0239}
\dcref{Corollary 7.7}
\group{proving-the-global-estimates}

If \((\mathcal M^n,g_0)\) is a Riemannian manifold such that \(|\Rm[g_0]|_{g_0}\le K\), then the unique Ricci flow solution \(g(t)\) with \(g(0)=g_0\) exists at least for \(t\in[0,c/K]\), where \(c>0\) depends only on \(n\).
\end{corollary}

\begin{proof}
\source{chow-knopf-ricci-flow:page-0239}
\sketch Apply Corollary~\dcref{Corollary 7.5} to the maximal solution starting at \(g_0\). The curvature bound can double only after time \(c/K\), so the solution cannot break down before that time.
\end{proof}

\section*{3. The Compactness Theorem}

\begin{theorem}[Compactness theorem for Ricci flow]
\label{ch07-thm-compactness-theorem}
\source{chow-knopf-ricci-flow:page-0244, chow-knopf-ricci-flow:page-0245}
\dcref{Theorem 7.8}
\group{the-compactness-theorem}

Let
\[
\{\mathcal M_i^n,g_i(t),O_i,F_i:i\in\mathbb N\}
\]
be a sequence of complete solutions to the Ricci flow existing for \(t\in(\alpha,\omega)\), where \(-\infty\le \alpha<0<\omega\le\infty\). Each solution is marked by an origin \(O_i\in \mathcal M_i^n\) and a frame \(F_i=\{e_1^i,\dots,e_n^i\}\) at \(O_i\) orthonormal with respect to \(g_i(0)\). Suppose that there exists \(K<\infty\) such that
\[
\sup_{\mathcal M_i^n\times(\alpha,\omega)}|\sect[g_i]|\le K
\]
for all \(i\), and suppose further that there exists \(\delta>0\) such that
\[
\inj_{g_i(0)}(O_i)\ge \delta
\]
for all \(i\). Then there exists a subsequence which converges in the pointed category to a complete solution
\[
\{\mathcal M_\infty^n,g_\infty(t),O_\infty,F_\infty\}
\]
of the Ricci flow existing for \(t\in(\alpha,\omega)\), with the same curvature and injectivity-radius bounds at the limit, and with convergence in every \(C^m\) norm on compact subsets after pullback by the associated diffeomorphisms.
\end{theorem}

\begin{proof}
\source{chow-knopf-ricci-flow:page-0244, chow-knopf-ricci-flow:page-0245}
\sketch The source states the compactness theorem and records its convergence conclusion, but defers the detailed proof to a successor volume. The statement is included here as part of the chapter's consequences of the derivative estimates.
\end{proof}
